\documentclass[UTF8,zihao=-4]{ctexart}
\usepackage[a4paper,margin=2.2cm]{geometry}
\usepackage{xcolor}
\usepackage{hyperref}
\usepackage{enumitem}
\usepackage{fontspec}
\usepackage{amsmath}
\IfFontExistsTF{Microsoft YaHei}{\setCJKmainfont{Microsoft YaHei}}{\setCJKmainfont{SimSun}[BoldFont=SimHei]}
\IfFontExistsTF{Microsoft YaHei UI}{\setCJKsansfont{Microsoft YaHei UI}}{}
\IfFontExistsTF{Consolas}{\setmonofont{Consolas}}{}
\definecolor{linkblue}{RGB}{30,80,145}
\hypersetup{colorlinks=true,linkcolor=linkblue,urlcolor=linkblue,pdfborder={0 0 0}}
\setlength{\parindent}{2em}
\setlength{\parskip}{0.25em}
\linespread{1.16}
\setlist{itemsep=0.2em,topsep=0.35em,leftmargin=2.5em}
\pagestyle{plain}

\title{04\_平滑函数与主定理}

\author{}

\date{}

\begin{document}

\maketitle

\section*{04 平滑函数与主定理}

\subsection*{题目原文}

了解平滑函数性质，应用主定理求解一般式：

\[T(n) = bT(n/c) + f(n)\]

\subsection*{考查类型}

概念题和递推分析题。

\subsection*{递推含义}

递推式：

\[T(n) = bT(n/c) + f(n)\]

表示一个规模为 $n$ 的问题被分成 \texttt{b} 个规模为 \texttt{n/c} 的子问题，合并或额外处理代价为 $f(n)$。

其中通常有：

\[b \ge  1\]
\[c > 1\]

关键比较对象是：

\[n^{\log_c b}\]

它表示递归树中叶子层的总规模量级。

\subsection*{主定理三种情形}

令：

\[\alpha = \log_c b\]

第一种：$f(n)$ 比 $n^\alpha$ 多项式意义下更小。

\[f(n) = O(n^{\alpha - \varepsilon})\]

其中 $\varepsilon > 0$，则：

\[T(n) = \Theta(n^\alpha)\]

第二种：$f(n)$ 与 $n^\alpha$ 同阶。

\[f(n) = \Theta(n^\alpha)\]

则：

\[T(n) = \Theta(n^\alpha \log n)\]

第三种：$f(n)$ 比 $n^\alpha$ 多项式意义下更大，并满足正则条件。

\[f(n) = \Omega(n^{\alpha + \varepsilon})\]

且存在常数 $q < 1$，使得：

\[b f(n/c) \le  q f(n)\]

则：

\[T(n) = \Theta(f(n))\]

\subsection*{平滑函数的作用}

在主定理相关证明中，平滑性质用于处理 $n$ 不是 $c$ 的整数幂时的情况，并保证函数在按常数比例缩放时增长不发生剧烈跳变。

常见的平滑函数包括：

\[n^{d}\]
\[n^{d} \log^{k} n\]

其中 \texttt{d} 和 \texttt{k} 是常数，且函数在讨论区间内非负。

这些函数满足按常数倍缩放时仍保持同一量级：

\[f(cn) = \Theta(f(n))\]

\subsection*{例子}

归并排序：

\[T(n) = 2T(n/2) + n\]

这里：

\[b = 2\]
\[c = 2\]
\[\alpha = \log_2 2 = 1\]
\[f(n) = n = \Theta(n^{1})\]

属于第二种情形：

\[T(n) = \Theta(n \log n)\]

二分查找：

\[T(n) = T(n/2) + 1\]

这里：

\[b = 1\]
\[c = 2\]
\[\alpha = \log_2 1 = 0\]
\[f(n) = 1 = \Theta(n^{0})\]

属于第二种情形：

\[T(n) = \Theta(\log n)\]

\subsection*{答题模板}

\begin{enumerate}
\item 写出 \texttt{b}、$c$、$f(n)$。
\item 计算 $\alpha = \log_c b$。
\item 比较 $f(n)$ 和 $n^\alpha$。
\item 判断主定理情形。
\item 写出 $T(n)$ 的 $\Theta$ 结果。
\end{enumerate}

\subsection*{例题}

\subsubsection*{例题 1}

题目：求解：

\[T(n) = 3T(n/2) + n\]

解答：

这里：

\[b = 3\]
\[c = 2\]
\[\alpha = \log_2 3\]
\[f(n) = n\]

因为：

\[n = O(n^{\log_2 3 - \varepsilon})\]

可取 $\varepsilon = \log_2 3 - 1$。因此属于主定理第一种情形：

\[T(n) = \Theta(n^{\log_2 3})\]

\subsubsection*{例题 2}

题目：求解：

\[T(n) = 4T(n/2) + n^{2}\]

解答：

这里：

\[b = 4\]
\[c = 2\]
\[\alpha = \log_2 4 = 2\]
\[f(n) = n^{2}\]

因为：

\[f(n) = \Theta(n^\alpha)\]

属于主定理第二种情形：

\[T(n) = \Theta(n^{2} \log n)\]

\subsubsection*{例题 3}

题目：求解：

\[T(n) = 2T(n/2) + n^{2}\]

解答：

这里：

\[b = 2\]
\[c = 2\]
\[\alpha = \log_2 2 = 1\]
\[f(n) = n^{2}\]

$f(n)$ 多项式意义下大于 $n^\alpha$。检查正则条件：

\[2f(n/2) = 2(n/2)^{2} = n^{2}/2 \le  (1/2)f(n)\]

取 $q = 1/2$，正则条件成立。属于主定理第三种情形：

\[T(n) = \Theta(n^{2})\]

\subsection*{常见误区}

第三种情形不能只比较大小，还要检查正则条件。

第二种情形要求同阶。若出现 $n^\alpha \log^k n$，需要按课堂给出的扩展主定理处理。

\end{document}
